\documentclass[12pt]{amsart}
\usepackage[letterpaper,margin=1in]{geometry}
%%%%%%%%%%%%%%%%%%%%%%%%%%%%%%%%%%%%%%%%%%%%%%%%%%%%%%%%%%%%%%%%%%%%%%%%%%%%
\usepackage{listings}
% \lstdefinestyle{pystyle}{
%        language=Python,
%        basicstyle=\ttfamily,
%        keywordstyle=\color{blue},
%        stringstyle=\color{red},
%        commentstyle=\color{green},
%        numbers=left,
%        numberstyle=\tiny\color{gray},
%        stepnumber=1,
%        showstringspaces=false
%    }
%%%%%%%%%%%%%%%%%%%%%%%%%%%%%%%%%%%%%%%%%%%%%%%%%%%%%%%%%%%%%%%%%%%%%%%%%%%%

\usepackage{amsmath,amssymb,amsfonts}
\usepackage{bm}
\usepackage{bbm}
\usepackage[config,labelfont={rm,rm},textfont=rm]{caption,subfig}

\usepackage{amsmath,amssymb,amsthm,latexsym,mathrsfs}

\usepackage{color,graphicx,colortbl}
\usepackage{bm,bbm}

\usepackage[latin1]{inputenc}
\usepackage[english]{babel}

\theoremstyle{definition}
\newtheorem{theorem}{Theorem}[section]
\newtheorem{lemma}[theorem]{Lemma}
\newtheorem{problem}[theorem]{Problem}
\newtheorem{proposition}[theorem]{Proposition}
\newtheorem{corollary}[theorem]{Corollary}
\newtheorem{definition}[theorem]{Definition}
\newtheorem{algorithm}[theorem]{Algorithm}
\newtheorem{example}[theorem]{Example}
\newtheorem{remark}[theorem]{Remark}
\newtheorem{exercise}[theorem]{Exercise}
\newtheorem{observation}[theorem]{{\color{red}Observation}}

\renewcommand{\labelenumi}{\textbf{\arabic{enumi}.}}
\renewcommand{\labelenumii}{\textbf{\arabic{enumi}\alph{enumii}.}}
\everymath{\displaystyle}

\title{An idea for solution. Question(s) for Zhilan \& John}
\author{\ldots \and Yuliya \and Ludmil}

\begin{document}

\maketitle

\begin{abstract}
We present a possible algorithm for minimizing the spectral radius of a class of non-negative matrices appearing in a math bio application. 
\end{abstract}

\section{Preliminaries}

We recall here the Perron-Frobenius theorem about non-negative and irreducible matrices because this is an essential result to be used. The proof of this famous theorem is found at many places, for example, check Varga~\cite[Theorem 2.7]{2000Varga-aa}. Please, pay particular attention to item 3. in the statement of the theorem. 
\begin{theorem}[Perron-Frobenius]\label{pf}
  Let $\bm A\geq 0$ be a non-negative and irreducible matrix. Then
  \begin{enumerate}
  \item $\bm A$ has positive real eigenvalue equal to $\rho(\bm A)$.
  \item The corresponding eigenvector has nonnegative elements (positive in fact).
  \item $\rho(\bm A)$ increases when any entry of $\bm A$ increases.
  \item $\rho(\bm A)$ is a simple (not repeating) eigenvalue of $\bm A$.
  \end{enumerate}
\end{theorem}

{\color{magenta} Similar result when the matrix is only nonnegative and might be reducible. We state two results here and refer the reader to~\cite[Section 8.3]{1985Horn_Johnson}
}
\begin{theorem}\label{reducible}
If $A\in \mathbb{R}^{n\times n}$ is a nonnegative, then $\rho(A)$ is an eigenvalue of A and there 
is a nonnegative nonzero vector $x$ such that $Ax=\rho(A)x$.
\end{theorem}
\begin{proof}\cite[Theorem 8.3.1]{1985Horn_Johnson}
Let $\mathbbm{1}$ be the column vector of all ones. Then we consider the matrix 
\[
A_{\varepsilon}(\varepsilon)=A+\varepsilon \mathbbm{1}\mathbbm{1}^*, \quad \varepsilon>0.
\]
The matrix $A_{\varepsilon}$ is a full matrix and therefore irreducible satisfying all assumptions of the Perron-Frobenius Theorem~\ref{pf}. Taking the limit when $\varepsilon\to 0_+$ then shows the result (the eigenvalues of a matrix are continuously depending on the matrix entries). 
\end{proof}

Regarding the other items in Theorem~\ref{pf} when compared to Theorem~\ref{reducible} is that Theorem~\ref{pf}(4.) may fail to hold; while Theorem~\ref{pf}(3.) holds but "increasing" should be replaced with "non-decreasing".

Next Corollary (found in~\cite[Corollary~8.3.3]{1985Horn_Johnson}) could be useful for computations:

\begin{corollary}
If $A\in\mathbb{R}^{n\times n}$ is nonnegative, then
\begin{equation}\label{max-min}
\rho(A) = \max_{\substack{x\ge 0 \\ x\neq 0}} 
\min_{\substack{1 \leq i  \leq n \\ x_i\neq 0}} \frac{(A x)_i}{x_i}.
\end{equation}
\end{corollary}
\section{Constraints}
We also have $3m$ constraints. Let $S_{ij}(T-1)$, $V_{ij}(T-1)$, and 
$P_{ij}(T-1)$ denote the numbers of individuals in the susceptible $(S)$, vaccinated ($V$), and $(P)$ classes {\color{red} at the end of week $T-1$} (which are the solutions of the ODEs by using the $\chi$ values in the week $T-1$). We need to find optimal $\chi$ values for week $T$. Then the constraints are:
\begin{equation}\label{eq:c}
\begin{aligned}
&\sum_{j=1}^n \chi_{1ij} S_{ij}(T-1) = b_{1i}, \quad i=1:m, \\
&\sum_{j=1}^n \chi_{2ij} V_{ij}(T-1) = b_{2i}, \quad i=1:m, \\
&\sum_{j=1}^n \chi_{3ij} P_{ij}(T-1) = b_{3i}, \quad i=1:m, 
\end{aligned}
\end{equation}
where $b_{ki}$ $(k=1, 2, 3)$ are the total available vaccine doses of type $k$ at location $i$ for week $T$. To unify the notation, let us denote 
\[
s_{1,ij}=S_{ij}(T-1), \quad 
s_{2,ij}=V_{ij}(T-1),\quad
s_{3,ij}=P_{ij}(T-1).
\]
The constraints then take the form:
\begin{equation}\label{eq:c-1}
\begin{aligned}
&\sum_{j=1}^n s_{k,ij}\chi_{k,ij} = b_{k,i}, \quad i=1:m, \quad k=1:3.
\end{aligned}
\end{equation}

We convert $(k,ij)$ to one index, $\ell=(k-1)nm+(i-1)n+j$, $k=1:3$, $i=1:m$, $j=1:n$, i.e. 
$l=1, \ldots, 3nm$ and denote the new variables with $\widehat{\bm{\xi}}$ and $\widehat{\bm{\sigma}}$.

\begin{equation}\label{eq:c-1}
\begin{aligned}
&(I_m\otimes  \widehat{\bm\sigma}_k)\widehat\xi_{(k-1)mn+1:kmn} = 
b_{(k-1)m+1:km},\quad (\widehat{\bm{\sigma}}_1)_{(i-1)n+j:in+j}=\sigma_{1,ij}.
\end{aligned}
\end{equation}
Furthermore, we set
\[
\begin{aligned}
&\chi_{k,ij}\rightarrow\widehat\xi_{(k-1)nm+n(i-1)+j}=\widehat\xi_{\ell},\\
&\sigma_{k,ij}\rightarrow\widehat C_{(k-1)nm+n(i-1)+j}=\widehat C_{\ell},\\
& k=1:3,\quad i=1:m,\quad j=1:n,\quad \mbox{and}\quad \ell=1:3nm,
\end{aligned}
\]
and then the constraints take the form:
\begin{equation*}
\begin{aligned}
&\sum_{\ell=(k-1)nm+(i-1)n+1}^{(k-1)nm+in} \widehat C_{\ell} \widehat\xi_{\ell}=b_{(k-1)m+i}, \quad i=1:m,\quad k=1:3. 
\end{aligned}
\end{equation*}
Let us change the variables $\xi_\ell=\frac{\widehat C_{\ell}}{b_{(k-1)m+i}}\widehat\xi_{\ell}$, for $i=1:m$ and $\ell=(k-1)nm+(i-1)n+j$, \quad $k=1:3$
\begin{equation*}
\begin{aligned}
&\sum_{\ell=(k-1)nm+(i-1)n+1}^{(k-1)nm+in} \xi_{\ell}=1, \quad i=1:m,\quad k=1:3. 
\end{aligned}
\end{equation*}
In short, the constraints can be written as ({\color{red} check!}
\begin{equation*}
\begin{aligned}
(I_m\otimes\mathbbm{1}_n\otimes I_3)\bm{\xi}=\mathbbm{1}_{3nm}. 
\end{aligned}
\end{equation*}



Notice also that the following identity holds:
\begin{equation}
    \text{NGM}=D(\bm{\chi})A_c, \quad\mbox{where}\quad \bm{\chi}=
\begin{pmatrix}
    \bm{\chi}_1\\
    \bm{\chi}_2\\
    \bm{\chi}_3    
\end{pmatrix}.
\end{equation}
The matrix $A_c$, is independent of the variables $\bm{\chi}_k$, $k=1:3$.
Below $n_c$ is the number of rows in $A_c$.
\begin{enumerate}
    % \item Similarly for the new variables $\bm{\xi}_k$. 
    % We can identify $\bm{\xi}_k$ by changing the indexing to one index, by rows. 
    \item Moreover, we have that the rows in $A_c$ repeat, that is every "odd" numbered row, equals the next "even" numbered row. 
    \item We can reduce $\text{NGM}$ by removing the kernel. First, the kernel is the range of the matrix $P_0$ defined as
    \begin{equation}
        P_0=
        \frac{1}{\sqrt{2}}\begin{pmatrix}
        1\\-1
        \end{pmatrix}\otimes I_{p},\quad
        P=\frac{1}{\sqrt{2}}\begin{pmatrix}
        1\\1
        \end{pmatrix}\otimes I_{p},\quad p=n_c/2 
    \end{equation}
    Clearly, 
    \begin{equation}
        A_cP_0=0, \quad P_0^{\prime}P=0, 
        \quad P'P=P_0^{\prime}P_0=I_{p}.     
    \end{equation}
    \item The eigenpair $(\lambda,\bm{\varphi})$ with $\lambda \neq 0$ then 
    satisfies: 
    \[
    D(\bm{\chi})A_c \bm{\varphi}=\lambda \bm \varphi.
    \]
    Since $\bm{\varphi}\in[\operatorname{Range}(A_c)]=[\operatorname{Ker}(A_c^T)]^\perp$ this is equivalent to the following smaller eigenproblem:
    \begin{equation}
        D_p(\bm{\chi}) A_p\bm{x} =\lambda \bm{x}, \quad D_p=P'DP, \quad A_p=P'A_cP, \quad \bm{x}\in \mathbb{R}^p. 
    \end{equation}
\end{enumerate}
\section{A special case}
We now focus on the problem when $\bm{\chi}_2=\bm{\chi}_3=0$.
\begin{enumerate}
    \item A straightforward computation (denoting $\bm{\chi}=\bm{\chi}_1$.) shows that:
    \begin{equation}
        (D_p)_j=(1-\varepsilon_2)+\varepsilon_2 f(\bm{\chi}_j), \quad f(t)=\frac{\omega_v}{\omega_v+t}.
    \end{equation}
    \item We observe that $1-\varepsilon_2 < f(t) < 1$ for all $t>0$.
    \item Our goal is to minimize $\rho$, the spectral radius of $D_pA_p$, which is the same as the the spectral radius of $DA_c$. As mentioned earlier. This is also the spectral radius of $A_pD_p$ or the spectral radius of $D_p^{1/2}A_pD_p^{1/2}$.
    \begin{equation}\label{max-min-1}
\rho = \max_{\substack{x\ge 0 \\ x\neq 0}} \min_{\substack{1 \leq i  \leq n \\ x_i\neq 0}} \frac{( D_p A_px)_i}{x_i}=
\max_{\substack{x\ge 0 \\ x\neq 0}} \min_{\substack{1 \leq i  \leq n \\ x_i\neq 0}} \frac{( A_pD_px)_i}{x_i}.
\end{equation}
\item For $i=1:p$, the following relations are obvious:
    \begin{equation}
    \begin{aligned}
        \min_{1\le i\le p} \frac{(D_p A_px)_i}{x_i}&= \min_{1\le i\le p}\left[(1-\varepsilon_2)\frac{(A_px)_i}{x_i}+
        \varepsilon_2 f(\bm{\chi}_i)\frac{(A_px)_i}{x_i}\right].
    \end{aligned}        
    \end{equation}
\item Clearly, taking the max with respect to $x$ we obtain that
    \begin{equation}
    \begin{aligned}
        \rho(\bm{\chi})&=\max_{x\ge 0}
        \min_{1\le i\le p} \frac{(D_p A_px)_i}{x_i}\\
        &= \max_{x\ge 0}\min_{1\le i\le p}\left[(1-\varepsilon_2)\frac{(A_px)_i}{x_i}+
        \varepsilon_2\frac{f(\bm{\chi}_i)(A_px)_i}{x_i}\right]\\
        &=\max_{x\ge 0}\min_{1\le i\le p}\left[(1-\varepsilon_2)+
        \varepsilon_2f(\bm{\chi}_i)\right]\frac{(A_px)_i}{x_i}\\
        &=\max_{x\ge 0}\min_{1\le i\le p}\left[1-\varepsilon_2\left(1-f(\bm{\chi}_i)\right)\right]\frac{(A_px)_i}{x_i}\\
        &=\max_{x\ge 0}\min_{1\le i\le p}\left[1-\varepsilon_2\frac{\bm{\chi}_i}{\omega_v+\bm{\chi}}_i\right]\frac{(A_px)_i}{x_i}
    \end{aligned}        
    \end{equation}
    % \item The above may not be the best approach. Let us state a conjecture: For any $x$ we can find $\bm{\chi}$ such that 
    % \(\left(1-\frac{\bm{\chi}_i}{\bm{\chi}_i+\omega_v}\right)\) and $(Ax)_i$, $i=1:p$ are at minimum at one and the same index. But then we don't know what to do with it...
    \item Consider $g(t)=1-\varepsilon_2\left(\frac{t}{t+\omega_v}\right)$; we have
    \begin{equation}
    \begin{aligned}
        g(t) &  = 1-\varepsilon_2\left(1-\frac{\omega_v}{t+\omega_v}\right)=1-\varepsilon_2+\frac{\varepsilon_2\omega_v}{t+\omega_v}, \\
        g'(t) & =-\frac{\varepsilon_2\omega_v}{(t+\omega_v)^2}<0, \quad\mbox{for all}\quad t\ge 0.
        \end{aligned}
    \end{equation}

\end{enumerate}

\section{Nonsense}
\begin{enumerate}
\item Let us continue the discussion replacing $D_pA_p$ with $A_pD_p$ in the definition of the spectral radius~\eqref{max-min-1}. Let $\operatorname{diag}$ and its right inverse $\operatorname{Diag}$ be defined as:
    \begin{equation}
        \begin{aligned}
        &\mbox{If}\quad X=\sum_j d_j e_j e_j^*+\sum_{i\neq_j}x_{ij}e_ie_j^*, 
            \quad\mbox{then}\quad \operatorname{diag}(X):=(d_1,\ldots,d_n)^*\\            
        &\mbox{If}\quad \bm{d}=(d_1,\ldots,d_n)^* 
            \quad\mbox{then} \quad \operatorname{Diag}(\bm{d}):=\sum_j d_j e_j e_j^*.
        \end{aligned}        
    \end{equation}
    \item Let us change the notation a bit and set 
    \[
    G(\bm{\chi})=(g(\chi_1),\ldots,g(\chi_n)),\qquad
    D_G(\bm{\chi})=\operatorname{Diag}(G(\bm{\xi}).
   \]
   With this notation, noting that $\operatorname{Diag}(a)b=\operatorname{Diag}(b)a$ we have
    \begin{equation}
        \begin{aligned}
            &A_p D_p x = A_pD_G(\bm{\chi}) x = A_p\operatorname{Diag}(x)G(\bm{\chi})=A_pD_xG(\bm{\chi}),\\
            &\frac{[A_pD_xG(\bm{\chi})]_i}{x_i}=\left[ D_x^{-1}A_pD_xG(\bm{\chi})\right]_i
        \end{aligned}
    \end{equation}
\end{enumerate}

\section{Lagrange multipliers}
For a given $x$ , let us denote $Z_x=D_x^{-1}A_cD_x$ and the rows of this matrix are $z_{j}$, $j=1:n$. We consider the following $n$ minimization problems (with constraints): 
\begin{equation}\label{22}
    \begin{aligned}
        &\bm{\widehat\chi}_j=\arg\min_{\bm{\xi}}f_j(\bm{\xi}),\quad f_j(\bm\xi):=\langle z_j,G(\bm{\xi})\rangle, \\
        &\mbox{subject to}\quad W^T \bm{\xi}-b=0, \quad -\bm{\xi}\le 0.
    \end{aligned}
\end{equation}

\section{KKT conditions}
Let us drop the index $j$ in the equation~\eqref{22} above so we can have a simpler notation. We have,
\begin{equation}
    \begin{aligned}
        &\bm \chi=\arg\min_{\bm{\xi}}f(\bm{\xi}),\quad f(\bm\xi):=\langle z,G(\bm{\xi})\rangle, \\
                &\mbox{subject to}\quad W^T \bm{\xi}-b=0, \quad -\bm{\xi}\le 0.
    \end{aligned}
\end{equation}
The Lagrangian associated with this problem is:
\begin{equation}\label{lagrange-0}
\mathcal{L}(\bm{\xi},\bm{\mu},\bm{\lambda}):=
\left\langle z, G(\bm{\xi})\right\rangle+
\left\langle \bm \lambda, W^T\bm{\xi}-\bm{b}\right\rangle-
\left\langle \bm \mu, \bm{\xi}\right\rangle
\end{equation}
We also have the conditions
\begin{equation}\label{cond-lagrange}
    \bm\mu \ge 0; \qquad \langle\bm\mu,\bm \xi\rangle =0 ,\quad \bm{\xi}\ge 0.
\end{equation}

\begin{equation}
    \frac{\partial G(\bm{\xi})}{\partial \xi_m} = -\delta_{mk}\left[\frac{\omega_v}{(\xi_m+\omega_v)^2}\right] 
\end{equation}

Differentiating the objective function then gives:
\begin{equation}\label{grad-o}
\begin{aligned}
    \frac{\partial }{\partial \xi_m}f(\bm{\xi})&=\sum_{k} z_{k}\frac{\partial  }{\partial \xi_m}g(\xi_k)
    =-\sum_{k} z_{k}\delta_{km}\frac{\omega_v}{(\xi_m+\omega_v)^2}
    =-z_{m}\frac{\omega_v}{(\xi_m+\omega_v)^2}
\end{aligned}
\end{equation}
For the second derivatives we have
\begin{equation}\label{hessian-o}
\begin{aligned}
    \frac{\partial^2 }{\partial \xi_m\xi_l}f(\bm{\xi})&=
    -    \frac{\partial }{\partial \xi_l}\left(z_{m}\frac{\omega_v}{(\xi_m+\omega_v)^2}\right)
    =\operatorname{Diag}\left(\frac{z_l\omega_v}{(\xi_l+\omega_v)^3}\right).
\end{aligned}
\end{equation}
This shows that the Hessian is a positive definite {\color{red}(could it be semi-definite?)}. That is then as follows:
\begin{equation}
\begin{aligned}
    &\frac{\partial }{\partial \xi_m}f(\bm{\xi},\bm{\lambda},\bm{\mu})=
    (W\bm{\lambda})_m-z_{m}\frac{\omega_v}{(\xi_m+\omega_v)^2}-\mu_m
    =0, \quad m=1:n\\
    &W^T\bm{\xi}-b =0, \quad
    \bm\mu \ge 0; \quad \langle\bm\mu,\bm \xi\rangle =0 ,\quad \bm{\xi}\ge 0.
\end{aligned}
\end{equation}
We thus need to solve the following system of equations+inequalities
\begin{equation}\label{conds-0}
\begin{aligned}
    &(W\bm{\lambda})_m-z_{m}\frac{\omega_v}{(\xi_m+\omega_v)^2}-\mu_m=0, \quad m=1:n\\
    &W^T\bm{\xi}-b =0, \quad
    \bm\mu \ge 0; \quad \langle\bm\mu,\bm \xi\rangle =0 ,\quad \bm{\xi}\ge 0.
\end{aligned}
\end{equation}

We multiply by $\xi_m$ and divide by $\mu_m^2$ the first equation in~\eqref{conds-0} and we use that $\xi_m\mu_m=0$. Let $\mathfrak{l}$ and $\mathfrak{m}$ 
be the sets where $\xi_l>0$ for $l\in \mathfrak{l}$ and $\mu_m>0$ for $m\in\mathfrak{m}$. We then have the following:
\begin{equation}\label{conds-10}
\begin{aligned}
    &\frac{1}{\mu_m^2}(W\bm{\lambda})_m-z_{m}\frac{\omega_v}{(\mu_m\xi_m+\mu_m\omega_v)^2}-\frac{\mu_m}{\mu_m^2}=0,\\
    &\frac{1}{\mu_m^2}(W\bm{\lambda})_m-z_{m}\frac{1}{\mu_m^2\omega_v}-\frac{1}{\mu_m}=0,\\
    &(W\bm{\lambda})_m-\frac{z_{m}}{\omega_v}-\mu_m=0,\qquad m\in\mathfrak{m}\\
    &\xi_l(W\bm{\lambda})_l-\xi_l z_{l}\frac{\omega_v}{(\xi_l+\omega_v)^2}=0, for l \in \mathfrak{l}.
\end{aligned}
\end{equation}

\begin{equation}\label{conds-1}
\begin{aligned}
    &\xi_l(W\bm{\lambda})_l-\xi_l z_{l}\frac{\omega_v}{(\xi_l+\omega_v)^2}-=0, \quad\mbox{for}\quad l \in \mathfrak{l},\\
    &(W\bm{\lambda})_m-\frac{z_{m}}{\omega_v}-\mu_m=0,\quad\mbox{for}\quad m\in\mathfrak{m}.
\end{aligned}
\end{equation}
Therefore, we have
\begin{equation}\label{conds-1}
\begin{aligned}
    &(W\bm{\lambda})_l-z_{l}\frac{\omega_v}{(\xi_l+\omega_v)^2}=0, \quad l\in \mathfrak{l},\\
    &(W\bm{\lambda})_m-\frac{z_m}{\omega_v}-\mu_m=0, \quad m\in \mathfrak{m}.
\end{aligned}
\end{equation}

% \begin{equation}\label{conds-1}
% \begin{aligned}
%     &(W\bm{\lambda})_l-z_l+z_{l}\left(1-\frac{\omega_v}{(\xi_l+\omega_v)^2}\right)&=0,
%     \quad &l\in \mathfrak{l},\\
%     &(W\bm{\lambda})_m-z_m-\mu_m&=0, \quad & m\in \mathfrak{m}.
% \end{aligned}
% \end{equation}

% \begin{equation}\label{conds-2}
% \begin{aligned}
%     &w_l + \frac{z_l\xi_l}{\xi_l+\omega_v}&=0,\quad &l\in \mathfrak{l},\\
%     &w_m - \mu_m&=0, \quad & m\in \mathfrak{m}.
% \end{aligned}
% \end{equation}

% \begin{equation}\label{conds-3}
% \begin{aligned}
%     \xi_l&=-\frac{w_l}{w_l+z_l}\omega_v,\quad &l\in \mathfrak{l},\\
%     \mu_m&=w_m, \quad & m\in \mathfrak{m}.
% \end{aligned}
% \end{equation}

% Using that $z_l>0$, the condition $\xi_l>0$ for $l\in \mathfrak{l}$ shows that $w_l$ and $(w_l+z_l)$ must have opposite signs. This is impossible if $w_l>0$ and is only possible if
% \begin{equation}\label{conds-4}
% \begin{aligned}
% &w_l<0, \quad\mbox{and}\quad |(W\bm{\lambda})_l-z_l|<z_l,\quad l\in\mathfrak{l},\\
% &-z_l<(W\bm{\lambda})_l-z_l<z_l,\quad \mbox{i.e.}\quad (W\bm{\lambda})_l<2\bm{z},\quad l\in\mathfrak{l},\\
% &w_l=(W\bm{\lambda})_l-z_l<0,\quad \mbox{i.e.}\quad (W\bm{\lambda})_l<z_l, \quad l\in\mathfrak{l}\\
% &\mbox{Hence},\quad 0<(W\bm{\lambda})_l<z_l, \quad l\in\mathfrak{l}.
% \end{aligned}
% \end{equation}
% Finally, using that for $m\in \mathfrak{m}$ we have $0<\mu_m<w_m=(W\bm{\lambda})_m-z_m$ we obtain that
% \begin{equation}\label{z1}
%     \begin{aligned}
%         &(W\bm{\lambda})_l<z_l, \quad l\in \mathfrak{l}\\
%         &z_m<(W\bm{\lambda})_m, \quad l\in \mathfrak{m}
%     \end{aligned}
% \end{equation}

From conditions~\eqref{conds-10} and changing $\xi_l\leftarrow\frac{\xi_l}{\omega_v}$ and 
$z_l\leftarrow\frac{z_l}{\omega_v}$ we find that
\begin{equation}
    \begin{aligned}
       ( W\bm\lambda)_l &= \frac{z_l}{(\xi_l+1)^2},\quad l\in \mathfrak{l},\quad\mbox{and}\\ 
       ( W\bm\lambda)_m &= z_m+\mu_m, \quad m\in \mathfrak{m}.
    \end{aligned}
\end{equation}





\section{Data details}
Each block $F_{ijpq}$ of the matrix $F$ represents the interactions between susceptible people of age $j$ at location $i$ with infected people of age $q$ at location $p$).

Let us fix $(i,j,p,q)$ and look at the definition of $\bm F_{ijpq}$:
\begin{equation}\label{f-ijpq-1}
\bm{F}_{ijpq}=
\begin{pmatrix}
0 & 0 &  Y_{ij}\Psi_{qI}\eta_{a}c_{ijpq} & \ldots & Y_{ij}\Psi_{qI}\eta_{m}c_{ijpq}\\
0 & 0 &  Z_{ij}\Psi_{qI}\eta_{a}c_{ijpq} & \ldots & Z_{ij}\Psi_{qI}\eta_{m}c_{ijpq}\\
0 & 0 &  0 & \ldots & 0\\
\vdots & \vdots &  \vdots & \ldots & \vdots\\
0 & 0 &  0 & \ldots & 0
\end{pmatrix}
\end{equation}
For the case of $m=4$ and $n=4$, the first row  of block matrices of  $F$  are 
\tiny{
\[
(F_{1111} \ \ F_{1112} \ \ F_{1113} \  \ F_{1114} \  \
F_{1121} \ \  F_{1122} \ \ F_{1123}  \ \ F_{1124}   \ \
F_{1131} \ \   F_{1132} \ \  F_{1133} \ \  F_{1134} \ \
F_{1141} \ \   F_{1142} \ \  F_{1143} \ \  F_{1144}).
\]
}
{\color{red}(ZF note: it is a row vector of block matrices {\bf each block} with dimension $8\times 8$. It describes the interactions between people of age $j=1$ at location $i=1$ with infected people of age $q=1,2,3,4$ at location $p=1,2,3,4$.)}
\tiny{
\begin{equation}\label{F-def-1}
\begin{pmatrix} 
0 & 0 & Y_{ij}\Psi_{qI}\eta_{a}c_{ijpq} & Y_{ij}\Psi_{qI}\eta_{s}c_{ijpq} &Y_{ij}\Psi_{qI}c_{ijpq} 
&Y_{ij}\Psi_{qI}\eta_{h}c_{ijpq}  &
Y_{ij}\Psi_{qI}\eta_{pm}c_{ijpq}  & Y_{ij}\Psi_{qI}\eta_{m}c_{ijpq}  
\end{pmatrix}.
\end{equation}
}
The values of $Y_{ij}$ and $Z_{ij}$ are computed as follows:
\begin{equation}\label{yz-ij-def}
\begin{aligned}
Y_{ij} &= \Psi_{jS}s_ia_{ij}\beta_j \frac{S_{ij}^*}{N_{ij}}, \quad i=1:m,\; j=1:n\\
Z_{ij} &= \Psi_{jS} s_i a_{ij}\beta_j \frac{(1-\varepsilon_{1})V_{ij}^*+(1-\varepsilon_{3})P_{ij}^*}{N_{ij}}.
\end{aligned}
\end{equation}
Here, we have denoted
\begin{equation}\label{si-def}
\begin{aligned}
\frac{S_{ij}^*}{N_{ij}}&= \left\lbrack
1+\frac{\chi_{1,ij}}{\omega_v+\varepsilon_2\chi_{2,ij}}
\left(1+\frac{\varepsilon_{2}\chi_{2,ij}}{\omega}+
\frac{\varepsilon_{2}\chi_{2,ij}}{\omega}\left\lbrack
1+\frac{\varepsilon_{3}\chi_{3,ij}}{\omega}
\right\rbrack
\right)
\right\rbrack^{-1}\\
\frac{V_{ij}^*}{N_{ij}}&=  \frac{\chi_{1,ij}}{\omega_{v}+\varepsilon_2\chi_{2,ij}} \frac{S_{ij}^*}{N_{ij}}\\
\frac{P_{ij}^*}{N_{ij}}&= \frac{\varepsilon_{2}\chi_{2,ij}}{\omega}\frac{V_{ij}^*}{N_{ij}}=
\frac{\varepsilon_{2}\chi_{2,ij}}{\omega}
\frac{\chi_{1,ij}}{\omega_{v}+\varepsilon_2\chi_{2,ij}} \frac{S_{ij}^*}{N_{ij}} \\
 \frac{R_{ij}^*}{N^*_{ij}}&=
\frac{\chi_{1ij}}{\omega_v+\varepsilon_2 \chi_{2ij}}
\frac{\varepsilon_2 \chi_{2ij}}{\omega}
\Big(1+\frac{\varepsilon_3 \chi_{3ij}}{\omega}\Big)
\frac{S_{ij}^*}{N^*_{ij}}.
\end{aligned}
\end{equation}


\begin{observation}
The way $F_{ijpq}$ are defined if we look at one $16\times 8$ block column for fixed $(p,q)$ where each block is $8\times 8$:
\begin{equation}
\begin{pmatrix}
F_{11pq}\\
F_{12pq}\\
\vdots\\
F_{14pq}\\
F_{21pq}\\
\vdots\\
F_{24pq}\\
\vdots\\
F_{44pq}
\end{pmatrix}
\end{equation}
The non-zero columns in the above $128\times 8$ matrix are proportional to each other. So this is a rank one matrix and it is seen from its definition. Since there are 16 such columns, this shows that the rank of $F$ is no more than $16$, because we have 16 such block columns. 
\end{observation}

\hrulefill

\centerline{To be continued}

\hrulefill


The full matrix $F$:
\tiny{
 \[
Row 1: (F_{1111} \ \ F_{1112} \ \ F_{1113} \  \ F_{1114} \  \
F_{1121} \ \  F_{1122} \ \ F_{1123}  \ \ F_{1124}   \ \
F_{1131} \ \   F_{1132} \ \  F_{1133} \ \  F_{1134} \ \
F_{1141} \ \   F_{1142} \ \  F_{1143} \ \  F_{1144})
\]
\[
Row 2: (F_{1211} \ \ F_{1212} \ \ F_{1213} \  \ F_{1214} \  \
F_{1221} \ \  F_{1222} \ \ F_{1223}  \ \ F_{1224}   \ \
F_{1231} \ \   F_{1232} \ \  F_{1233} \ \  F_{1234} \ \
F_{1241} \ \   F_{1242} \ \  F_{1243} \ \  F_{1244})
\]
\[
Row 3: (F_{1311} \ \ F_{1312} \ \ F_{1313} \  \ F_{1314} \  \
F_{1321} \ \  F_{1322} \ \ F_{1323}  \ \ F_{1324}   \ \
F_{1331} \ \   F_{1332} \ \  F_{1333} \ \  F_{1334} \ \
F_{1341} \ \   F_{1342} \ \  F_{1343} \ \  F_{1344})
\]

\quad \vdots
\[
Row 16: (F_{4411} \ \ F_{4412} \ \ F_{4413} \  \ F_{4414} \  \
F_{4421} \ \  F_{4422} \ \ F_{4423}  \ \ F_{4424}   \ \
F_{4431} \ \   F_{4432} \ \  F_{4433} \ \  F_{4434} \ \
F_{4441} \ \   F_{4442} \ \  F_{4443} \ \  F_{4444})
\]
}

{\color{magenta} There are questions in regard to this (maybe). 
\begin{itemize}
    \item 
The matrix \verb|mF| in the \verb|Mathematica| notebook is with dimensions $(4,4,4,4,8,8)$. One can use \verb|ArrayReshape[mF,{128,128}]| in \verb|Mathematica| to convert $F$ to $128 \times 128$. The resulting matrix is of rank $4$ only! 
\item If we do not convert the matrix $F$ with \verb|ArrayReshape[]|, then Mathematica cannot compute the eigenvalues, can it? 
\end{itemize}
}%% end magenta

\bibliographystyle{plain}
\bibliography{bib_mb}

%%%%%%%%%%%%%%%%%%%%%%%%%%%%%%%%%%%%%

\newpage
\section{Answered\label{sec:answers}}
\begin{itemize}
\item If $\bm V$ is defined as in the document you gave us then we have that $\bm V$ is non-zero for $i=p$, and zero for $i\neq p$ (but nothing is said about $j$ and $q$). The definition of $\bm V$  from the document suggests the following form of $\bm V$
 
\begin{equation}\label{V-1}
\bm V = (V_{ijpq}):=\begin{pmatrix}
V_{1111} & V_{1112} & V_{1113} &  0 &  0 & 0 \\
V_{1211} & V_{1212} & V_{1213} &  0 &  0 & 0 \\
V_{1311} & V_{1312} & V_{1313} &  0 &  0 & 0 \\
0 & 0 & 0 & V_{2121} &  V_{2122} & V_{2123} \\
0 &  0 & 0 & V_{2221} &  V_{2222} & V_{2223} \\
0 &  0 & 0 & V_{2321} &  V_{2322} & V_{2323}
\end{pmatrix}
\end{equation}
In such case, we are unsure about the positivity of $\bm V^{-1}$; such positivity is not guaranteed, which bothers us a bit. However, $\bm H$ can still be positive, but this will not be easy to establish, or, at least, it is not obvious to us.

\item On the other hand, if $\bm V$ is non-zero only when $i=p$ and $j=q$, then the positivity of $\bm{V}^{-1}$ can be guaranteed, and $\bm{V}$ is block diagonal.
\begin{equation}\label{V-2}
\bm V =  (V_{ijpq}) := \begin{pmatrix}
V_{1111} & 0 & 0 &  0 &  0 & 0 \\
0 & V_{1212} & 0 &  0 &  0 & 0 \\
0 & 0 & V_{1313} &  0 &  0 & 0 \\
0 & 0 & 0 & V_{2121} &  0 & 0 \\
0 &  0 & 0 & 0 &  V_{2222} & 0 \\
0 &  0 & 0 & 0 &  0 & V_{2323}
\end{pmatrix}
\end{equation}
\item {QUESTION: Which of the two definitions of $\bm V$ is the correct one (equation~\eqref{V-1} or equation~\eqref{V-2})}. 
\item {\bf ANSWER:} The $\bm V$ given in~\eqref{V-2} is the correct one. 

\item QUESTION(s) 1: What is the number of columns in $\bm{F}_{ijpq}$? Is it $8$ because of the two zero columns added? The $m$ in $\eta_{qm}$ is not the same as the number of blocks $m=2$ defined in the first item, right? Anyway, how is the number of rows and columns in $F_{ijpq}$ determined as a function of the sites and the age groups  ($m$ and $n$  from the first item above).

\item {\bf Answer:} $F_{ijpq}$ are the derivatives of the RHS functions of the model equations on page 1 with respect to the 8 variables in the vector \[
v_{ij}=( E_{ijs}, E_{ijp}, I_{ija}, I_{ijps}, I_{ijs}, I_{ijh}, I_{ijpm}, I_{ijm} ) 
\]
(see the top of page 2) evaluated at $v_{ij}=(0,0,\cdots, 0)$
and $S_{ij}=S^*_{ij}, V_{ij}=V^*_{ij}, P_{ij}=P^*_{ij}$.  The letter $m$ in subscripts of variables (e.g., $E$ or $I$) and parameters (e.g., $\eta$) are independent of the number of age groups ($n$) or sites ($m$).

\item Question 2: Does $\bm{V}$ depend on $\bm{\chi}$? It will be great to know either way. It will be easier to deal with the optimization problem (minimizing the entries of $\bm F$ if $\bm V$ is independent of $\bm \chi$, but this is not a show-stopper.

\item {\bf Answer:} The matrix $V$ does NOT depend on $\chi$.
 
\item Question 3: Which entries of $\bm F$ depend on $\bm \chi$. All? Some? How?

\item {\bf Answer:} Because all elements of $F$ depend on either $Y$ or $Z$, where $Y$ and $Z$ depend on $S^*/N=s^*$, or $V^*/N=v^*$, or $P^*/N=p^*$, which depend on $\chi$ as shown in line 5 on page 2, all non-zero elements of F depend on $\chi$.

\item {\bf Additional explanations:} In the attached pdf, I highlighted the items related to the v vector, $s^*$, $v^*$ etc. $Y_{ij}$ and $Z_{ij}$.
 
The constraint on $\chi$ is illustrated on page 4 (Equation (215)). But for the real model, the constraints consist of $m$ equations ($m$ is the number of sites) of the form:
\[
\chi_{i1}N_{i1} +\chi_{i2}N_{i2}+\ldots +\chi_{in}N_{in}=c_i,     \quad i=1:m.
\] 
Here, $N_{ij}$ is the total number of individuals of location  $i$ and age $j$. When you need these $N_{ij}$ values as well as other parameter values, we can ask John.
 
Although $Y$ and $Z$ do not depend on $(pq)$, the factor $c_{ijpq}$ in $F_{ij}$ depend of parameters associated with $p$ or $q$. Thees parameters $c_{ijpq}$ link populations of different ages and different sites.
 
\end{itemize}
\end{document}

% \section{Code}

% %\lstset{style=pystyle}

% %\begin{lstlisting}
% %\lstinputlisting[firstline=300,lastline=500]{Code/data_math_bio.py}
% \lstinputlisting[language=Python,
%        basicstyle=\ttfamily,
%        keywordstyle=\color{blue},
%        stringstyle=\color{red},
%        commentstyle=\color{green}]{Code/z_math_bio.py}
% %\end{lstlisting}
% \end{document}


% We now state a helpful result with an obvious proof. 
% \begin{lemma}\label{l-1}
% Let $\bm A\ge 0$, $\bm B\ge 0$, and $\bm C\ge 0$ be nonnegative matrices. If $\bm B$ and $\bm C$  satisfy the inequality $\bm B \ge \bm C$ (to be understood element-wise),  then 
% \begin{equation}
% \bm{A}\bm B\ge \bm{A}\bm C,\qquad\mbox{where the inequality holds element-wise.}
% \end{equation}
% \end{lemma} 
% \begin{proof} Since $(\bm B - \bm C)\ge 0$ and $\bm A\ge 0$, it is obvious that $\bm A(\bm B - \bm C)\ge 0$ which is what we wanted to show. 
% \end{proof}

% Here is a possible path to the solution of the problem in hand:
% \begin{enumerate}
% \item As discussed at NSF: If $(z,\bm\varphi)$ is eigenpair for $\bm F\bm V^{-1}$, then $(z,\bm \psi)$ is 
% an eigenpair for $\bm V^{-1}\bm F$ with $\bm \psi=\bm V\bm \varphi$. So the eigenvalues are the same and we will focus on minimizing 
% the dominant eigenvalue $z=\rho\left(\bm V^{-1} \bm F\right)$.  
% \item Let $\bm H(\bm \chi)=\bm{V}^{-1}\bm F(\bm \chi)$ be the matrix for which we would like to minimize the dominant eigenvalue with respect to $\bm \chi$. Since $\bm V^{-1}\ge 0$ (this is easy to prove), from Lemma~\ref{l-1} and the Perron-Frobenius Theorem (item 3.) we have the following result.
% \begin{theorem}
% Let  $\widehat{\bm \chi}$ and $\widetilde{\bm \chi}$ be two values of the parameters satisfying the constraints. We then have
% \begin{equation}\label{monotonicity} 
% \mbox{If} \quad 
% \bm F(\widehat{\bm\chi})\ge \bm F(\widetilde{\bm\chi})\quad\mbox{then}\quad
% \rho\left(\bm H(\widehat{\bm\chi})\right) \ge \rho\left(\bm H(\widetilde{\bm\chi})\right) 
% \end{equation}
% \end{theorem}
% \begin{proof}
% The proof follows from Lemma~\ref{l-1} and the Perron-Frobenius Theorem.
% \end{proof}
% \item Basically, we need to minimize the entries of $\bm F$ with respect to $\bm{\chi}$ satisfying the constraints. That is it.
% \item From this observation, it is clear that that {\bf we do not need to know the dominant eigenvalue or the gradient with respect to $\bm \chi$ to find the optimal $\bm \chi$}.
% \end{enumerate}
